\scanpage{162}\mspage{20}
% archive sheet number: 162
Ces conditions impliquent que les objets $0$-connexes \deleted{sont}\inserted{de} $\unsure{\mathscr X|\mathcal X}^{\wedge}$ (en particulier les éléments de $S$) sont $\infty$-connexes.

\uline{Corollaire} \unsure{2.7|2.8}. \inserted{(De l'équivalence $b\Longleftrightarrow d$)} \deleted{Soit une} Soit $X$ un objet de $\mathrm{Cat}$. Pour que $X$ soit filtrante à g. il faut et s. que $\mathrm{Top}(X)$ soit \deleted{tot.} totalement $0$-connexe et \deleted{tot.} localement \deleted{\unsure{tot.}}\inserted{irréductible} \unsure{.}

% slanted
\margin{Exemple de $X$ tot. $0$-connexe, mais non filtrante, i.e. non loc. irréductible ? La catégorie $\Delta$ des \unsure{simplexes}, tot. $0$-connexe, mais n'est pas filtrante car \diagram{\begin{tikzcd}[ampersand replacement=\&,column sep=large]
\Delta^0=e\arrow[r,shift left,"\delta_0"]\arrow[r,shift right,swap,"\delta_1"]\&\Delta^1
\end{tikzcd}} ne peut être \unsure{égalisé}. Donc $\Delta^\circ=I$ est telle que $I\longrightarrow I^E$ soit cofinal pour tt ens fini $E$, mais $I$ non filtrante.}

[ou encore que \inserted{$X$ soit tot. $0$-connexe et que} pour tt $x\in\mathrm{Ob}\,X$, la catégorie induite $X_{/x}$ soit \inserted{irréductible, cela est $=$ (tot. $0$-connexes)}\deleted{\unsure{connexe} et que} pour deux objets \deleted{$x$ et $y$ \unsure{connexes} \illegible{1}} \deleted{\unsure{deux} objets $a,b$ de $X$,} \unsure{ou encore} \inserted{$a,b\in\mathrm{Ob}\,X$} \unsure{que} $x\in\mathrm{Ob}\,X$, $a\times b$ \inserted{dans $X^{\wedge}$} soit $0$-connexe, et $a\times_x b$ (où $a\to x\leftarrow b$ dans $X$) soit $0$-connexe]

% slanted
\margin{\deleted{\uline{NB} les conditions entre [ ] (qui expriment que les $X_{/x}$ sont filtrantes et \unsure{$X$} $0$-connexe), n'impliquent pas que $X$ soit filtrante}}

\uline{Démonstration} de 2.6. On a les implications \deleted{du} $\pm$ tautologiques

\diagram{\begin{tikzcd}[ampersand replacement=\&,row sep=small,column sep=normal]
\&a\arrow[rr,Leftrightarrow]\arrow[dl,Rightarrow,swap,"\text{par }2.5"]\arrow[d,Rightarrow]\&\&a'\arrow[d,Rightarrow]\\
\unsure{b}\arrow[d,Rightarrow]\&\unsure{c|\varepsilon}\arrow[rr,Leftrightarrow,"(\Leftarrow\,\text{par }2.\deleted{4}\inserted{3}.)"]\arrow[dl,Rightarrow]\&\&c'\arrow[dlll,Rightarrow,"\deleted{\substack{\unsure{\uparrow\,?}\\\unsure{(\Leftarrow\,\text{par }1.)|(\Leftarrow\,\text{par }2.4.)}}}"]\\
\mathrm{PS}1\&\&\&
\end{tikzcd}}
\deleted{\unsure{$\Longleftrightarrow$}}

\deleted{\deleted{Notons que} Montrons quand même les deux implications $c')\Longrightarrow\mathrm{PS}1$ et $b\Longrightarrow\mathrm{PS}1$. Supposons c'), et soit un diagramme \diagram{\begin{tikzcd}[ampersand replacement=\&,row sep=small,column sep=small]
a\arrow[dr]\&\\
\&x\\
b\arrow[ur]\&
\end{tikzcd}} dans $S$, par c') $a\times_x b$ \inserted{dans $\mathscr X^{\wedge}$} est $0$-connexe, donc non vide, ce qui signifie, \unsure{soit|sait} (comme $S$ est génératrice et ses éléments sont ``non vides'') \deleted{\unsure{qu'il}} qu'il existe $c$ dans $S$ et $c\longrightarrow a\times_x b$, \deleted{\illegible{2}} i.e. \diagram{\begin{tikzcd}[ampersand replacement=\&,row sep=small,column sep=small]
\&a\arrow[dr]\&\\
c\arrow[ur,dashed]\arrow[dr,dashed]\&\&x\\
\&b\arrow[ur]\&
\end{tikzcd}} rendant le carré}
